%% =====================================================================
%%  T-16-01  Concession and Retreat
%%  -------------------------------------------------------------------
%%  One idea per file. Core is mandatory. Slots in canonical order.
%%  Max 700 words. No layout commands. No filler. New source material
%%  for this entry goes in sources/<BOOK>/T-16-01.tex -- never by
%%  rewriting this file (docs/RULES.md R0.1).
%% =====================================================================
%% @kind: topic
%% @id: T-16-01
%% @title: Concession and Retreat
%% @subtitle: A demand refused and then reduced gets more than the reduced demand alone
%% @chapter: c16
%% @order: 10
%% @status: stable
%% @sources: B4
%% @updated: 2026-09-19

\topic{T-16-01}{Concession and Retreat}{A demand refused and then reduced gets more than the reduced demand alone}

\begin{core}
Open with a demand you expect refused, then withdraw it in favour of a smaller one. The smaller demand is granted at a rate it would not have achieved on its own, because the withdrawal is read as a concession and a concession invites one.
\end{core}
\begin{mechanism}
Two principles combine and neither requires sincerity.

Reciprocity supplies the obligation: a move toward you creates a duty to move toward the other party, and reducing a demand is such a move whether or not the reduced demand was the intended one all along. Contrast supplies the calibration: the second demand is judged against the first rather than on its own, so it arrives already looking moderate.

The important property is that both respond to the \emph{form} of the move rather than to the intent behind it. Knowing the sequence is scripted does not switch either off, which is why the technique survives being recognised. What it purchases is not agreement with a position but a felt duty to meet a halfway point that was never halfway.
\end{mechanism}
\begin{tells}
  \item The opening demand is larger than the other party could have expected you to accept.
  \item It is withdrawn quickly and without visible reluctance.
  \item The replacement is presented as a compromise.
  \item You feel you gained something by refusing and cannot say what.
  \item Both demands serve the same interest, and the second was available from the start.
\end{tells}
\begin{moves}
  \item Open above your real target, and treat the opening as an instrument rather than as a position to be defended.
  \item Retreat in one clear step. Several small reductions read as ordinary negotiation and lose the concession effect entirely.
  \item Make the retreat visible and slightly costly-looking; a withdrawal that appears effortless invites no return.
  \item Ask for the return in a different currency from the concession, so no direct comparison is available.
\end{moves}
\begin{counters}
  \item Re-classify the move: ask whether the second demand is a concession or simply the offer, and evaluate it as though it had arrived first.
  \item Decline to reciprocate a retreat that costs the other party nothing. A reduction that leaves their interest intact was never a concession.
  \item Test it with a retreat of your own. Reduce your position slightly and see whether they reduce again; a conciliator will, a script will not.
  \item Slow the exchange down. The duty to answer a concession is strongest inside a single conversation.
  \item Ask for the reduced demand in writing, on its own, dated. Detached from the opener it usually loses most of its force.
\end{counters}
\begin{cost}
The opener is exposed if the first demand is absurd, because an obviously planted anchor destroys the credibility of everything that follows --- including the honest parts of the position. It also teaches the counterpart to read every later reduction as scripted, which removes its value in any negotiation that continues. Cheap once against a stranger; expensive repeatedly against someone you will meet again.
\end{cost}
\begin{field}
A subscription pitched annually, declined, and replaced with a monthly one closes at a rate the monthly offer would not have closed had it been made first.
\end{field}

\sources{B4}
\seesources
\seealso{T-14-02, T-19-01, T-24-01, T-16-02}
